% Chapter 2: Theoretical Foundations — Mid-Term Report Phase 3

\chapter{Theoretical Foundations}

\section{Verifiable Credentials}

\subsection{Concept of Verifiable Credentials}

A Verifiable Credential (VC) is a standardized data model specified by the World Wide Web Consortium (W3C) to represent a set of claims made by an issuer about a subject, enabling third-party verifiers to cryptographically authenticate those claims without direct real-time communication with the issuing authority~\cite{W3CVC}. Unlike legacy paper diplomas or unauthenticated scanned PDF files, Verifiable Credentials are cryptographic artifacts: digital signatures generated by the issuer are bound directly to the credential payload. Consequently, any verifier possessing the issuer's public key or registered decentralized identifier (DID) can autonomously verify the integrity and provenance of the attestation. The W3C published the official Verifiable Credentials Data Model v2.0 as a formal recommendation~\cite{W3CVC}.

A standard Verifiable Credential consists of three core components:
\begin{enumerate}
    \item \textbf{Credential Metadata}: Context parameters defining the issuer identifier, issuance timestamp, expiration date, schema definition, and credential type (e.g., \texttt{vct}).
    \item \textbf{Credential Subject}: A structured payload containing the assertions and attributes granted to the recipient (e.g., degree title, graduation date, student identifiers, academic achievements).
    \item \textbf{Proof}: A cryptographic signature or proof structure (such as an ECDSA signature or EIP-712 attestation) that seals the entire credential payload against unauthorized alteration.
\end{enumerate}

Canonical representations commonly adopt JSON-LD (JSON for Linked Data) or JSON Web Signature/Token (JWT) serialization formats.

\subsection{The Issuer--Holder--Verifier Ecosystem}

The W3C VC architecture formalizes a decentralized three-party trust triangle~\cite{W3CVC}:
\begin{itemize}
    \item \textbf{Issuer}: An authoritative organization (such as a university or educational institution) that asserts claims about a subject, constructs the credential, signs it with an authoritative cryptographic key, and delivers it to the holder.
    \item \textbf{Holder}: The subject (e.g., student or graduate) who receives and stores the credential in a digital wallet, retaining complete sovereignty over when, how, and with whom the credential is shared.
    \item \textbf{Verifier}: A relying party (e.g., an employer, academic admissions board, or background screening service) that receives a credential presentation from the holder, verifies the cryptographic signature against the issuer's public trust anchor, and validates its revocation status.
\end{itemize}

In the \textbf{BK Credential System}, this paradigm is realized natively: the university is registered on-chain as an authorized issuer, the student acts as the self-sovereign holder, and verifiers validate credentials through public web portals or CLI tools without requiring direct access to university internal databases.

\section{Selective Disclosure JSON Web Tokens (SD-JWT)}

\subsection{JWT Foundations and Privacy Over-disclosure Bottlenecks}

JSON Web Token (JWT) is an established open standard (RFC 7519) representing compact, URL-safe JSON claims secured by symmetric secrets or asymmetric key pairs. While JWT is ubiquitous in authentication and authorization infrastructures, utilizing raw JWTs for academic credentials introduces severe privacy risks: standard JWT payloads are entirely visible to any recipient. A holder presenting a JWT diploma cannot conceal extraneous personal attributes (such as date of birth, national identity numbers, or cumulative GPA) without invalidating the cryptographic signature over the payload, leading to widespread over-disclosure of Personally Identifiable Information (PII).

\subsection{The SD-JWT Mechanism}

Selective Disclosure for JWTs (SD-JWT) addresses this limitation by enabling holders to selectively disclose individual claims while preserving the mathematical validity of the issuer's original signature~\cite{RFC9901}. The core SD-JWT specification is standardized in IETF RFC 9901~\cite{RFC9901}, with its profile for Verifiable Credentials defined in \texttt{draft-ietf-oauth-sd-jwt-vc}~\cite{SDJWTVC}.

The operational mechanics of SD-JWT are structured as follows:
\begin{enumerate}
    \item For each private claim to be protected, the issuer generates a cryptographically secure 16-byte random salt.
    \item The issuer constructs a triple \texttt{[salt, claimName, claimValue]} and encodes it into a Base64URL string:
    \begin{equation}
    \text{Disclosure} = \text{base64url}(\text{JSON}([\text{salt}, \text{claimName}, \text{claimValue}]))
    \end{equation}
    \item The issuer computes the disclosure commitment digest:
    \begin{equation}
    \text{commitment} = \text{base64url}(\text{SHA-256}(\text{Disclosure}))
    \end{equation}
    and inserts it into the \texttt{\_sd[]} array within the signed JWT payload.
    \item All individual disclosures are appended to the signed JWT token, separated by tilde (\texttt{\textasciitilde{}}) delimiter characters.
\end{enumerate}

When presenting the credential, the holder strips undisclosed disclosures, transmitting only the base JWT and the subset of disclosures chosen for release. The verifier recomputes SHA-256 hashes of the supplied disclosures and asserts that they match the digests in the signed \texttt{\_sd[]} array.

\section{Blockchain and Smart Contracts}

\subsection{Blockchain Foundations and Immutability}

A blockchain is a decentralized, cryptographically linked distributed ledger where transactions are organized into sequential blocks and validated through consensus algorithms~\cite{Cachin2017Blockchain,Narayanan2016BitcoinBook}. Each block header incorporates the cryptographic hash of the preceding block along with a Merkle root representing all included transactions. Modifying historical records requires rewriting all subsequent blocks across a distributed majority of nodes, making retrospective falsification computationally and economically infeasible. This structural immutability makes blockchain an ideal foundation for hosting decentralized \textbf{Trust Anchors} --- tamper-proof public registries that provide cryptographic provenance without centralized single points of failure.

\subsection{Ethereum, EVM, and Smart Contracts}

Ethereum~\cite{EthereumWhitepaper} extends distributed ledger capabilities by integrating a Turing-complete decentralized execution engine known as the Ethereum Virtual Machine (EVM)~\cite{EthereumYellowPaper}. Smart contracts are autonomous, immutable programs deployed directly to the blockchain state~\cite{Szabo1997SmartContracts}. Once deployed, their bytecode executes deterministically across all validating network nodes, producing transparent and non-repudiable state transitions.

In the \textbf{BK Credential System}, smart contracts act as the decentralized Trust Anchor:
\begin{itemize}
    \item \texttt{IssuerRegistry.sol}: Governs authorized signing addresses belonging to accredited educational institutions.
    \item \texttt{CredentialRegistryV3.sol}: Persists Merkle root commitments corresponding to graduating batches and maintains bitmap revocation states.
\end{itemize}

Verifiers query these smart contracts directly via standard JSON-RPC endpoints to confirm issuer validity and credential status without relying on university internal servers.

\subsection{Sepolia Testnet Environment}

Sepolia is an established Ethereum test network utilizing a Proof-of-Stake consensus mechanism identical to the Ethereum Mainnet while using testnet Ether devoid of real-world monetary value. Deploying contracts to Sepolia enables empirical profiling of gas costs, transaction lifecycle latency, and state transitions under real-world EVM network conditions.

\section{Merkle Trees and Cryptographic Proofs}

\subsection{Merkle Tree Data Structure}

A Merkle tree is a binary tree data structure where every leaf node represents the cryptographic hash of a data element, and every non-leaf node represents the cryptographic hash of its concatenated child nodes~\cite{Narayanan2016BitcoinBook}. The single hash at the top of the tree, denoted as the \textbf{Merkle Root}, serves as a compact, immutable cryptographic digest representing the entire set of elements. Any modification to a single leaf irrevocably alters the resulting Merkle Root.

\subsection{Merkle Proofs and Membership Verification}

Merkle trees enable proving that a specific element belongs to a set of size $n$ without disclosing the remaining elements. A \textbf{Merkle Proof} (or inclusion proof) consists of the sequence of sibling hashes along the direct path from the target leaf node to the root.

A verifier verifies membership by:
\begin{enumerate}
    \item Computing the leaf hash of the candidate element.
    \item Iteratively concatenating and hashing the result with each sibling hash in the proof path.
    \item Comparing the calculated top hash against the trusted Merkle Root stored on-chain.
\end{enumerate}

The computational and space complexity of Merkle inclusion verification is strictly $\mathcal{O}(\log n)$, allowing efficient proof generation and verification for batches comprising tens of thousands of credentials.

\subsection{Application in Batch Credential Issuance}

The \textbf{BK Credential System} utilizes Merkle trees to aggregate thousands of individual graduate credentials into a single cohort batch. The issuer computes the leaf hash for each credential, constructs a Merkle tree using \texttt{@openzeppelin/merkle-tree}, and submits a single transaction anchoring the Merkle Root onto \texttt{CredentialRegistryV3}. Each graduate receives their credential bundle alongside an individualized Merkle proof. Verifiers validate batch membership locally against the on-chain Merkle Root, eliminating the need to store thousands of individual certificate hashes on the blockchain.

\section{EIP-712: Typed Structured Data Hashing and Signing}

\subsection{Motivation: Transitioning from SD-JWT to EIP-712}

Phase 2 utilized SD-JWT (Selective Disclosure JSON Web Token) with ECDSA over the secp256k1 curve (ES256K). However, ES256K is not universally standardized across standard OAuth/JWT libraries (where RFC 7518 specifies ES256 over NIST P-256), requiring custom library implementations.

EIP-712~\cite{EIP712} provides a standardized hashing and signing protocol for typed structured data within the Ethereum ecosystem. Its primary advantages include:

\begin{enumerate}
    \item \textbf{Native Web3 Integration}: Direct support in standard Ethereum wallets (e.g., MetaMask, WalletConnect) via \texttt{eth\_signTypedData\_v4}.
    \item \textbf{Battle-tested Reliability}: Built-in support in standardized libraries such as \texttt{ethers.js} v6 via \texttt{signTypedData} and \texttt{verifyTypedData}.
    \item \textbf{Type-Safe Schemas}: Explicit type definitions prevent field ordering ambiguity and type casting vulnerabilities.
    \item \textbf{Replay Attack Mitigation}: Binding signatures to a unique Domain Separator (\texttt{chainId} + \texttt{verifyingContract}) prevents signature reuse across networks and contracts.
\end{enumerate}

\subsection{Domain Separator and Type Hash}

EIP-712 uses a \textbf{Domain Separator} to uniquely scope signatures to a specific decentralized application, network, and target contract address:

\begin{lstlisting}
domainSeparator = keccak256(
    keccak256("EIP712Domain(string name,string version,
               uint256 chainId,address verifyingContract)")
    || keccak256("BKCredential")
    || keccak256("3")
    || uint256(chainId)            // 11155111 for Sepolia
    || address(verifyingContract)  // CredentialRegistryV3 address
)
\end{lstlisting}

Each data type $T$ is hashed into an immutable \textbf{Type Hash}:

\begin{lstlisting}
typeHash = keccak256(encodeType(T))
\end{lstlisting}

where \texttt{encodeType} represents the canonical string encoding of the struct and all referenced subordinate structs sorted alphabetically.

\subsection{Nested Struct Encoding}

Phase 3 adopts a \textbf{nested struct} design: \texttt{BkCredential} contains an embedded \texttt{PublicClaims} struct. This provides several benefits:

\begin{enumerate}
    \item \textbf{Type Safety}: Clearly identifies public claims as structured data rather than opaque byte hashes.
    \item \textbf{Extensibility}: New public claims can be added to \texttt{PublicClaims} without altering the top-level credential schema.
    \item \textbf{Standard Hashing}: Supported natively by EIP-712 recursive \texttt{hashStruct} rules.
\end{enumerate}

The struct hashing is formally computed as:

\begin{lstlisting}
hashStruct(publicClaims) = keccak256(
    typeHash(PublicClaims)
    || keccak256(bytes(vct))
    || keccak256(bytes(degreeTitle))
    || keccak256(bytes(graduationDate))
    || keccak256(bytes(honors))
)
\end{lstlisting}

The final EIP-712 signature hash is given by:

\begin{lstlisting}
signHash = keccak256("\x19\x01" || domainSeparator 
           || hashStruct(credential))
\end{lstlisting}

\section{Selective Disclosure on EIP-712}

\subsection{Salted-Hash Commitment Scheme}

Phase 3 replaces SD-JWT disclosures with a lightweight \textbf{salted-hash commitment} scheme embedded directly in the EIP-712 payload:

\begin{lstlisting}
commitment = keccak256(abi.encodePacked(salt, key, value))
\end{lstlisting}

Where:
\begin{itemize}
    \item \texttt{salt}: A cryptographically secure 32-byte pseudo-random value generated via \texttt{crypto.getRandomValues()} or \texttt{ethers.randomBytes(32)}, providing 256 bits of entropy to resist brute-force dictionary attacks.
    \item \texttt{key}: UTF-8 string representing the attribute name (e.g., \texttt{"fullName"}, \texttt{"studentId"}, \texttt{"dob"}).
    \item \texttt{value}: UTF-8 string representing the attribute value.
\end{itemize}

The \texttt{privateClaims[]} array in \texttt{BkCredential} stores all 32-byte commitment hashes. When a holder chooses to reveal a claim, they provide the triple \texttt{\{salt, key, value\}} (termed a \textbf{disclosure}). The verifier recomputes the commitment hash and confirms its membership in \texttt{privateClaims[]}.

\textbf{Security Guarantee}: Because \texttt{privateClaims[]} is embedded inside the signed EIP-712 message, any tampering with commitments invalidates the signature. Verifiers cannot derive undisclosed values without knowing the corresponding random 256-bit \texttt{salt}.

\subsection{Comparison with SD-JWT Disclosure}

Table~\ref{tab:sd-comparison} summarizes the improvements over the Phase 2 SD-JWT implementation.

\begin{table}[H]
\centering
\renewcommand{\arraystretch}{1.3}
\begin{tabular}{|p{4.2cm}|p{5.2cm}|p{5.2cm}|}
\hline
\textbf{Feature} & \textbf{Phase 2 (SD-JWT)} & \textbf{Phase 3 (Salted-Hash)} \\
\hline
Hash Algorithm & SHA-256 & keccak256 (EVM native) \\
\hline
Disclosure Format & Base64url(JSON([salt, key, value])) & \texttt{\{salt, key, value\}} plaintext object \\
\hline
Commitment Location & \texttt{\_sd[]} array in JWT payload & \texttt{privateClaims[]} in EIP-712 struct \\
\hline
Required Library & Custom SD-JWT Parser & \texttt{ethers.solidityPackedKeccak256} \\
\hline
Envelope Overhead & \texttt{\_sd\_alg}, JWT headers, formatting & Zero additional envelope overhead \\
\hline
Verification Complexity & Multi-step JWT parse & 1 line of keccak256 computation \\
\hline
\end{tabular}
\caption{Selective Disclosure Comparison: SD-JWT vs. Salted-Hash}
\label{tab:sd-comparison}
\end{table}

\section{Bitmap-Based Revocation}

\subsection{On-Chain Bitmap Mechanism}

In Phase 2, revocation status was stored via individual mapping entries: \texttt{mapping(bytes32 credId => bool revoked)}, consuming 1 cold storage slot (32 bytes = 256 bits) for every single revoked credential.

Phase 3 introduces an \textbf{on-chain bitmap revocation} scheme where each 256-bit \texttt{uint256} storage word tracks the status of 256 individual credentials:

\begin{lstlisting}
wordIndex = credentialIndex / 256
bitMask   = 1 << (credentialIndex % 256)

// Revoke credential at index:
_revocationBitmap[merkleRoot][wordIndex] |= bitMask;

// Check revocation status:
isRevoked = (_revocationBitmap[merkleRoot][wordIndex] 
             & bitMask) != 0;
\end{lstlisting}

For a graduating batch of 10,000 students, revocation storage requires only $\lceil 10,000 / 256 \rceil = 40$ words instead of 10,000 individual storage slots.

\subsection{Gas and Storage Optimization}

\begin{table}[H]
\centering
\renewcommand{\arraystretch}{1.3}
\begin{tabular}{|p{4.2cm}|p{5.2cm}|p{5.2cm}|}
\hline
\textbf{Metric} & \textbf{Phase 2 (Per-credential mapping)} & \textbf{Phase 3 (Bitmap)} \\
\hline
Storage per Credential & 1 slot (32 bytes) & 1/256 slot ($\sim$0.125 bytes) \\
\hline
Gas per Revocation (Single) & $\sim$45,000 gas (Cold SSTORE) & $\sim$5,000 gas (Warm SSTORE, bitwise OR) \\
\hline
Batch Revocation & $N \times 45,000$ gas & $\sim$90,000 gas total (4 items in same word) \\
\hline
Efficiency Gain & Reference baseline & $\sim$64$\times$ storage, $\sim$9$\times$ gas reduction \\
\hline
\end{tabular}
\caption{Revocation Storage and Gas Cost Comparison}
\label{tab:revocation-comparison}
\end{table}

\section{OpenZeppelin StandardMerkleTree}

\subsection{Double keccak256 Leaf Hashing}

Phase 3 integrates \texttt{@openzeppelin/merkle-tree} (StandardMerkleTree)~\cite{OZMerkleTree} to standardize tree generation and prevent cryptographic vulnerabilities such as second-preimage attacks through double hashing:

\begin{lstlisting}
leaf = keccak256(keccak256(abi.encode(value1, value2, ...)))
\end{lstlisting}

The initial \texttt{abi.encode} enforces unambiguous byte serialization, while the outer hash prevents intermediate nodes from colliding with valid leaf hashes.

\subsection{Comparison with Phase 2 Merkle Scheme}

\begin{table}[H]
\centering
\renewcommand{\arraystretch}{1.3}
\begin{tabular}{|p{4.2cm}|p{5.2cm}|p{5.2cm}|}
\hline
\textbf{Dimension} & \textbf{Phase 2} & \textbf{Phase 3} \\
\hline
Implementation & Custom TypeScript Tree & OpenZeppelin StandardMerkleTree \\
\hline
Leaf Serialization & \texttt{keccak256(credId || pcvCid)} & \texttt{keccak256(keccak256(abi.encode( credId, pubHash, privHash)))} \\
\hline
Double Hashing & No & Yes (Prevents 2nd preimage attacks) \\
\hline
Audit Status & Unaudited & Fully audited by OpenZeppelin \\
\hline
Integration Standard & Custom logic & Standardized npm ecosystem package \\
\hline
\end{tabular}
\caption{Merkle Tree Scheme Comparison}
\label{tab:merkle-comparison}
\end{table}